% ============================================================
% Chapter 1: Introduction
%   Section: Problem Statement
% Included from chapters/01-introduction/introduction.tex
% ============================================================

\section{Problem Statement}
\label{sec:problem}

Given a set of topics chosen by an editor, the system has to do what a
journalist of a constructive publication does: find the recent
articles on those topics that are worth offering, decide which of them
to analyse, and verify the claims each one makes before it is
published. Four properties of the domain make this harder than the
problem automated fact-checking usually studies.

\begin{enumerate}
    \item \textbf{Claims mix facts with interpretations.} A
    constructive article reports a study, a project, a figure or an
    event, and interprets it (``a hopeful step for the sector''). The
    first kind can be checked; the second cannot, and a system that
    tries to check it either wastes effort or produces verdicts with no
    meaning. The checkable claims have to be told apart from the rest
    before any evidence is sought.
    \item \textbf{Evidence is on the live web, often in one place.} The
    claims are recent and specific, so no fixed corpus holds their
    evidence: it has to be found on the open web at the time of
    checking. Frequently the only source is the press release of the
    organisation the story is about, which is not independent
    confirmation.
    \item \textbf{Two languages.} The sources publish in English and in
    Spanish, so extraction, claim selection, search and verification
    must work in both, and be measured in both.
    \item \textbf{False claims are rare.} Constructive outlets seldom
    publish a claim that the evidence refutes; errors are more often in
    a detail than in the substance. A system can look accurate by
    answering ``true'' to everything, and a benchmark from this domain
    has almost no refuted class to test it against.
\end{enumerate}

Formally, for each claim $c$ selected from an article published at time
$t_c$, the system retrieves a set of evidence documents $E$ from the
open web at time $t_v \geq t_c$ and predicts a verdict
$y \in \{\textsc{true}, \textsc{partially true}, \textsc{misleading},
\textsc{false}, \textsc{unverified}\}$ together with the subset of $E$
it rests on. \textsc{unverified} is a legitimate answer whenever the
evidence does not decide the claim, and it is preferred to a guess.
The gap between $t_c$ and $t_v$ matters: evidence published after the
claim was made, including a fact-check of the claim itself, can make a
claim look easier to verify than it was when it was published
\cite{glockner2022missing}.
